% !TeX root = ../../tesis.tex

% =============================================================================
%  B.1 — CONTRATOS DE VALIDACIÓN PYDANTIC
% =============================================================================
\section{Contratos de Validación de Datos (Pydantic)}
\label{anx:vft-validacion}

El contrato de datos de VFTModel está implementado en
\path{src/api/schemas/schemas.py} mediante cuatro enumeraciones cerradas y un
esquema principal de validación (\texttt{PropertiesSchema}). El diseño garantiza
que ningún dato proveniente del servidor Go llegue al constructor del grafo sin
haber pasado por la taxonomía del modelo. El Cuadro~
\ref{tab:vft-campos-validados} presenta los campos validados; el código
siguiente muestra las enumeraciones y el validador dinámico que imputa los
valores faltantes.

\begin{table}[H]
    \centering
    \caption[Campos validados en el contrato de datos de VFTModel]{Campos
    validados en el contrato de datos de VFTModel.
    \fuentefigura{\texttt{src/api/schemas/schemas.py} (VFTModel, 2026).}}
    \label{tab:vft-campos-validados}
    \begin{tabularx}{\textwidth}{|l|l|X|}
        \hline
        \textbf{Campo} & \textbf{Tipo} & \textbf{Descripción} \\
        \hline
        \texttt{sistema}                  & Enum (12 valores) &
            Identifica el sistema de transporte \\
        \texttt{tipo\_entidad}            & Enum (2 valores) &
            \texttt{estacion} o \texttt{ruta} \\
        \texttt{jerarquia\_transporte}    & Enum (4 niveles) &
            Jerarquía operativa del sistema \\
        \texttt{derecho\_de\_via}         & Enum (4 valores) &
            Exclusividad del carril de circulación \\
        \texttt{velocidad\_promedio\_kmh} & Flotante &
            Velocidad comercial en km/h \\
        \texttt{frecuencia\_minutos}      & Flotante &
            Intervalo promedio entre servicios \\
        \texttt{sentido}                  & Entero &
            Dirección de la ruta: ida (1), regreso (0) \\
        \texttt{es\_cetram}               & Booleano &
            Si la estación es un nodo de transferencia \textsc{cetram} \\
        \hline
    \end{tabularx}
\end{table}

\subsection*{Enumeraciones de taxonomía}

\begin{lstlisting}[style=estiloPython,
    caption={Enumeraciones cerradas del contrato de datos VFTModel.
             \texttt{src/api/schemas/schemas.py}.},
    label=lst:vft-enums]
class JerarquiaTransporte(str, Enum):
    masivo_pesado              = "masivo_pesado"
    masivo_mediano             = "masivo_mediano"
    superficie_convencional    = "superficie_convencional"
    superficie_baja            = "superficie_baja"

class DerechoVia(str, Enum):
    exclusivo  = "exclusivo"   # Metro, Suburbano, Cablebus
    confinado  = "confinado"   # Metrobus, Mexibus
    compartido = "compartido"  # Trolebus
    mixto      = "mixto"       # RTP, CC, Pumabus

class TipoEntidad(str, Enum):
    estacion = "estacion"
    ruta     = "ruta"

class SistemaTransporte(str, Enum):
    metro      = "METRO";    metrobus  = "MB"
    tren_lig   = "TL";       trolebus  = "TROLE"
    rtp        = "RTP";      cablebus  = "CBB"
    cc         = "CC";       mexicable = "MEXICABLE"
    mexibus    = "MEXIBUS";  suburbano = "SUB"
    interurb   = "INTERURBANO";  pumabus = "PUMABUS"
\end{lstlisting}

\subsection*{Validador dinámico y motor de imputación}

El \texttt{model\_validator} de Pydantic actúa como segunda línea de defensa:
detecta y repara campos nulos antes de que el dato llegue al constructor del
grafo. Las imputaciones cubren cuatro casos frecuentes en el \textsc{geojson} de
Apimetro.

\begin{lstlisting}[style=estiloPython,
    caption={Validador dinámico que repara datos incompletos antes de construir el
             grafo. \texttt{src/api/schemas/schemas.py}.},
    label=lst:vft-validator]
@model_validator(mode='after')
def validar_y_reparar_entidades(self):
    # 1. Nombres nulos
    if self.nombre is None:
        self.nombre = "Desconocido"
    # 2. Jerarquia huerfana
    if self.jerarquia_transporte is None:
        self.jerarquia_transporte = JerarquiaTransporte.superficie_convencional
    # 3. Imputacion para RUTAS (aristas fantasma)
    if self.tipo_entidad == TipoEntidad.ruta:
        if self.capacidad_vehiculo is None:
            self.capacidad_vehiculo = 60
        if self.derecho_de_via is None:
            self.derecho_de_via = DerechoVia.mixto
    # 4. Imputacion para ESTACIONES (nodos)
    elif self.tipo_entidad == TipoEntidad.estacion:
        if not self.alcaldia_municipio or self.alcaldia_municipio == "0.0":
            self.alcaldia_municipio = "Desconocida"
        if not self.tipo_estacion or self.tipo_estacion == "0.0":
            self.tipo_estacion = "Superficie"
    return self
\end{lstlisting}
